% ============================================================
\section{Conclusion}
% ============================================================

What follows is not a morality play of bad technology and good law. It is a claim about failed closure. Once output-only review cannot stabilize same-case-same-action, unresolved questions of authority and responsibility return to institutions. The Cascade provides the diagnostic vocabulary for those institutions to ask \textit{which layer failed} before they answer \textit{who bears responsibility}. Technical systems remain part of the problem, but they cannot by themselves close the questions of authority and responsibility that their failures open.

To move beyond model-centric regulation is therefore not to deny that models matter. It is to refuse to collapse the entire governance stack into the model-user dyad. Many failures that appear to be failures of model behavior are better understood as Override (a higher layer suppressing a lower), Inadequate Process Model (contamination or coverage gaps leave controllers with an incorrect or empty model of the hazard), or Compression (multiple layers combining to produce behavior no single layer intended). The Cascade distinguishes these patterns. Whether the institutional drift toward upstream responsibility is ethically sound lies outside this diagnostic frame---but the drift itself is visible, structural, and accelerating.

